\begin{description}
\item[(a)] Aldebaran is far enough away that the light rays from it can be
  considered to be parallel, thus the problem can be modelled by the following
  diagram:

  % FIGURE NEEDED: diagram of parallel light rays from Aldebaran grazing the
  % Moon (centre M, points C, D on the rays) and reaching the Earth (points B
  % = Bologna and P at latitude phi_1, declination delta_A, radius r);
  % 2023 theory solutions, page 26 (problem T10).

  where the light ray through $C$ passes through the centre of the Earth and
  thus $\delta_{\mathrm{A}}$ is the declination of Aldebaran; the ray through
  $D$ and $B$ (Bologna) represents the situation observed by Copernicus, with
  Aldebaran behind the Moon and $0.32'$ above the south edge; and the ray
  through $M$ and $P$ represents the situation where Aldebaran is exactly
  behind the centre of the Moon. $\varphi_{\mathrm{B}}$ is the geocentric
  latitude of Bologna, $\varphi_1$ is the latitude of the place $P$ and $r$ is
  the radius of the Earth (assume the Earth is spherical).

  Calculating the latitude $\varphi_1$ is thus a matter of simple geometry,
  where:
  \begin{align*}
    |CD| &= r\sin(\varphi_{\mathrm{B}} - \delta_{\mathrm{A}}),\\
    |CM| &= r\sin(\varphi_1 - \delta_{\mathrm{A}}),
  \end{align*}
  and the distance $|DM|$ is the fraction of the Moon's radius given by:
  \[
    |DM| = 1\,737\ \mathrm{km} \times
      \left(1 - \frac{0.32'}{31.5'/2}\right) = 1\,702\ \mathrm{km}.
  \]
  Therefore, since
  \begin{align*}
    |CM| &= |CD| + |DM| = 3\,099\ \mathrm{km} + 1\,702\ \mathrm{km}
          = 4\,801\ \mathrm{km},\\
    \Longrightarrow\ \varphi_1 &= \arcsin\left(\frac{4\,801}{6\,378}\right)
      + \delta_{\mathrm{A}} = 64.19^\circ.
  \end{align*}
  \hfill(6 points)

\item[(b)] Given the assumptions, the duration of the occultation will just
  depend on the difference between the tangential speeds of the observer,
  $v_{\mathrm{obs}}$, and Moon, $v_{\mathrm{Moon}}$, and on the diameter of
  the Moon.

  For the observer, the tangential speed is given by the circumference of
  their path at latitude $\varphi_1$ divided by the sidereal day:
  \[
    v_{\mathrm{obs}} = \frac{2\pi r_{\mathrm{Earth}} \cos\varphi_1}
                            {23^{\mathrm{h}}\,56\mathrm{m}\,4^{\mathrm{s}}}
                     = 729\ \mathrm{km/h}
  \]

  % FIGURE NEEDED: sketch of Aldebaran, the Moon (M, velocity v_Moon) and the
  % rotating Earth with observer P (velocity v_obs) at latitude phi_1;
  % 2023 theory solutions, page 27 (problem T10).

  Similarly for the Moon, using the semi-major axis of the Moon's orbit and
  the length of the sidereal month (from the table of constants):
  \[
    v_{\mathrm{Moon}} = \frac{2\pi \times 3.844 \times 10^5\ \mathrm{km}}
                             {27.321661\ \mathrm{d} \times 24}
                      = 3\,683\ \mathrm{km/h}
  \]
  The difference of speeds is then:
  \[
    v_{\mathrm{rel}} = v_{\mathrm{Moon}} - v_{\mathrm{obs}}
                     = 2\,954\ \mathrm{km/h}
  \]
  and the time taken to move across the lunar diameter
  ($2 \times 1737$\,km) at that speed is:
  \[
    t = \frac{2 \times 1\,737\ \mathrm{km}}{2\,954\ \mathrm{km/h}}
      = 1.176 \approx 1.18\ \mathrm{hours}.
  \]
  \hfill(5 points)

\item[(c)] To determine the topocentric angular velocity, we need to find the
  distance $d$ between the observer (at place $P$ at latitude $\varphi_1$) and
  the Moon:

  % FIGURE NEEDED: triangle Moon (M) - observer (P) - centre of the Earth,
  % with the angles rho, delta, delta', alpha, sides a, d, r and latitude
  % phi_1 marked; 2023 theory solutions, page 27 (problem T10).

  where $a$ is the semi-major axis of the Moon's orbit, $\delta$ is the
  geocentric declination of the Moon, and $\delta'$ is the topocentric
  declination of the Moon as seen by the observer. We can rewrite
  $\varrho = \delta - \delta'$ and
  $\alpha = 180^\circ - (\varphi_1 - \delta) - \varrho
          = 180^\circ - (\varphi_1 - \delta) - (\delta - \delta')
          = 180^\circ - (\varphi_1 - \delta')$,
  then using the sine theorem:
  \[
    \frac{d}{\sin(\varphi_1 - \delta)}
      = \frac{a}{\sin\left(180^\circ - (\varphi_1 - \delta')\right)}
      = \frac{r}{\sin(\delta - \delta')}
  \]
  Simplifying and substituting $\delta' = \delta_{\mathrm{A}}$:
  \[
    \frac{d}{\sin(\varphi_1 - \delta)}
      = \frac{a}{\sin(\varphi_1 - \delta_{\mathrm{A}})}
      = \frac{r}{\sin(\delta - \delta_{\mathrm{A}})}
  \]
  We can now calculate $\delta$:
  \[
    \delta = \arcsin\left(
        \frac{r\sin(\varphi_1 - \delta_{\mathrm{A}})}{a}\right)
      + \delta_{\mathrm{A}} = 16.08^\circ
  \]
  and $d$:
  \[
    d = a\,\frac{\sin(\varphi_1 - \delta)}
                {\sin(\varphi_1 - \delta_{\mathrm{A}})}
      = 380\,203\ \mathrm{km}.
  \]
  Thus the angular velocity is:
  \[
    \omega = \left(\frac{2\,R_{\mathrm{Moon}}}{d}\right)
             \left(\frac{1}{1.176\ \mathrm{h}}\right) \mathrm{rad/h}
           = \left(\frac{2 \times 1\,737}{380\,203}\right)
             \left(\frac{1}{1.176\ \mathrm{h}}\right)
             \left(\frac{180 \times 60}{\pi}\right)
           = 26.71 \approx 27\ \mathrm{arcmin/h}.
  \]
  \hfill(7 points)

\item[(d)] The result in part (c), 27\,arcmin/h, was calculated for the
  situation when the Moon and observer are moving in the same direction,
  reducing their relative velocity, and thus represents the lower limit of the
  range of topocentric angular velocities. The upper limit is given by the
  situation when the Moon and observer are moving in the opposite direction,
  i.e.\ when the Moon is due north.

  Using the calculation from part (b), we find that the relative linear speed
  is:
  \[
    v'_{\mathrm{rel}} = v_{\mathrm{Moon}} + v_{\mathrm{obs}}
                      = 3\,683 + 729 = 4\,412\ \mathrm{km/h}
  \]
  Neglecting the effect of the slightly greater distance between the observer
  and the Moon when the Moon is on the northern side of the sky, the angular
  speed will be greater by the ratio of the linear speeds,
  \[
    \omega' = \omega\,\frac{4\,416}{2\,954} = 39.89
            \approx 40\ \mathrm{arcmin/h},
    % sic: the official solution writes 4416 here; the relative speed derived
    % above is 4412 km/h.
  \]
  thus the range is approximately 27--40 arcmin/h. \hfill(3 points)

  If the additional distance is taken into account, the upper limit will be
  around 1\% less.

  Expressing the velocities as vectors,
  \[
    \vec{v}_{\mathrm{rel}} = \vec{v}_{\mathrm{Moon}} - \vec{v}_{\mathrm{obs}},
  \]
  the magnitude of the relative velocity,
  $v_{\mathrm{rel}} = |\vec{v}_{\mathrm{rel}}|$ is given by the scalar
  projection of $\vec{v}_{\mathrm{obs}}$ on $\vec{v}_{\mathrm{Moon}}$:
  \[
    v_{\mathrm{rel}} = v_{\mathrm{Moon}} - v_{\mathrm{obs}}\cos\alpha
  \]
  where $\alpha$ is the angle between the vectors. From this it can be seen
  that the relative velocity must be minimum for $\alpha = 0^\circ$ and
  maximum for $\alpha = 180^\circ$. These conditions are met on the local
  meridian, due south or due north, respectively. \hfill(4 points)
\end{description}
